%-------------------------
% Resume
% Author : Геннадий Воронин
%------------------------

\documentclass[letterpaper,11pt]{article}

\usepackage[T2A]{fontenc}
\usepackage[utf8]{inputenc}
\usepackage{latexsym}
\usepackage[empty]{fullpage}
\usepackage{titlesec}
\usepackage{marvosym}
\usepackage[usenames,dvipsnames]{xcolor}
\usepackage{verbatim}
\usepackage{enumitem}
\usepackage[hidelinks]{hyperref}
\usepackage{fancyhdr}
\usepackage[russian,english]{babel}
\usepackage{tabularx}
\usepackage{iftex}

\ifPDFTeX
  \input{glyphtounicode}
  \pdfgentounicode=1
\fi

\definecolor{cvblue}{HTML}{4493f8}
\newcommand{\resumelink}[2]{\href{#1}{\textcolor{cvblue}{#2}}}
\hypersetup{
  pdftitle={Геннадий Воронин — Resume (Data Engineer)},
  pdfauthor={Геннадий Воронин},
  pdfsubject={Data Engineering, resume},
  pdfkeywords={data engineering, resume}
}

\pagestyle{fancy}
\fancyhf{}
\fancyfoot{}
\renewcommand{\headrulewidth}{0pt}
\renewcommand{\footrulewidth}{0pt}
\setlength{\footskip}{8pt}

\addtolength{\oddsidemargin}{-0.5in}
\addtolength{\evensidemargin}{-0.5in}
\addtolength{\textwidth}{1in}
\addtolength{\topmargin}{-.5in}
\addtolength{\textheight}{1.0in}

\urlstyle{same}
\raggedbottom
\raggedright
\setlength{\tabcolsep}{0in}

\titleformat{\section}{
  \vspace{-4pt}\scshape\raggedright\large
}{}{0em}{}[\color{black}\titlerule \vspace{-5pt}]

\newcommand{\resumeItem}[1]{
  \item\small{{#1}}
}

\newcommand{\resumeSubheading}[4]{
  \vspace{-2pt}\item
    \begin{tabular*}{\linewidth}[t]{l@{\extracolsep{\fill}}r}
      \textbf{#1} & #2 \\
      {\small#3} & {\small #4} \\
    \end{tabular*}\vspace{-7pt}
}

\newcommand{\resumeProjectHeading}[2]{
  \vspace{-2pt}\item
    \begin{tabular*}{\linewidth}[t]{l@{\extracolsep{\fill}}r}
      \small#1 & \small #2 \\
    \end{tabular*}\vspace{-2pt}
}

\renewcommand\labelitemii{$\vcenter{\hbox{\tiny$\bullet$}}$}

\newcommand{\resumeSubHeadingListStart}{\begin{itemize}[leftmargin=0pt, label={}]}
\newcommand{\resumeSubHeadingListEnd}{\end{itemize}}
\newcommand{\resumeItemListStart}{\begin{itemize}[leftmargin=1.4em, topsep=1pt, partopsep=0pt, parsep=0pt, itemsep=0pt]}
\newcommand{\resumeItemListEnd}{\end{itemize}\vspace{-6pt}}
\newcommand{\resumeSimpleListStart}{\begin{itemize}[leftmargin=1.4em, topsep=0pt, partopsep=0pt, parsep=0pt, itemsep=-2pt]}
\newcommand{\resumeSimpleListEnd}{\end{itemize}}

%-------------------------------------------
\begin{document}

%----------HEADING----------
\begin{center}
  {\Large Геннадий Воронин} \\ \vspace{2pt}
  {\large Data Engineer} \\ \vspace{1pt}
  {\small
    \resumelink{https://t.me/unavailable_domain}{Telegram} $|$
    \resumelink{mailto:gndvrn@ya.ru}{gndvrn@ya.ru} $|$
    \resumelink{https://github.com/gndvrn}{GitHub}
  }
\end{center}

%-----------PROJECT EXPERIENCE-----------
\section{Опыт работы и проекты}
\resumeSubHeadingListStart

\resumeProjectHeading
{\resumelink{https://github.com/gndvrn/RESTAPI-Data-Batch-Pipeline}{\textbf{REST API Data Batch Pipeline}} $|$ Python, SQL, Airflow, dbt, ClickHouse, Docker, Superset}{проект}
\resumeItemListStart
\resumeItem{Реализовал учебный batch data pipeline: ingest из публичных REST API $\rightarrow$ ClickHouse $\rightarrow$ dbt $\rightarrow$ Superset}
\resumeItem{Выстроил слои staging $\rightarrow$ core $\rightarrow$ marts; построил аналитические витрины в dbt}
\resumeItem{Настроил оркестрацию в Airflow (2 DAG: extract/load и dbt run); покрыл ключевые поля dbt-тестами для контроля качества данных}
\resumeItem{Развернул multi-service stack в Docker Compose}
\resumeItemListEnd

\resumeProjectHeading
{\resumelink{https://github.com/gndvrn/e-commerce-dwh}{\textbf{E-commerce DWH}} $|$ Python, SQL, Airflow, ClickHouse, dbt, Docker, BPMN}{проект}
\resumeItemListStart
\resumeItem{Учебный проект --- хранилище e-commerce: спроектировал BPMN-модель компании; на её основе --- схему DWH (20+ таблиц) с четырьмя операционными источниками}
\resumeItem{Разработал генератор синтетических транзакций для имитации работы магазина; реализовал пакетный ETL: выгрузка в объектное хранилище, загрузка в ClickHouse, витрины в dbt, визуализация в Superset}
\resumeItemListEnd

\resumeProjectHeading
{\resumelink{https://github.com/gndvrn/docker-hadoop-spark-hive}{\textbf{Hadoop-кластер в Docker}} $|$ Hadoop, Hive, Spark, Docker, Bash}{проект}
\resumeItemListStart
\resumeItem{Развернул 11-сервисный Hadoop-кластер в Docker для курсового проекта: HDFS, YARN, Hive, Spark, JupyterLab.}
\resumeItem{Настроил Hive с metastore на PostgreSQL и Spark on YARN; автоматизировал запуск и обслуживание кластера bash-CLI ($\sim$500 строк).}
\resumeItemListEnd

\resumeProjectHeading
{\resumelink{https://github.com/gndvrn/kafka-spark-postgres-grafana-streaming}{\textbf{Streaming Data Pipeline}} $|$ Kafka, Spark, PostgreSQL, Grafana, Docker}{проект}
\resumeItemListStart
\resumeItem{Учебный проект --- доработал потоковый пайплайн мониторинга операционных данных интернет-магазина: Kafka $\rightarrow$ Spark $\rightarrow$ PostgreSQL $\rightarrow$ Grafana}
\resumeItem{Добавил доменную модель и SQL-схему; развернул сервисы в Docker Compose, подготовил дашборд в Grafana}
\resumeItemListEnd

\resumeSubHeadingListEnd

%-----------TECHNICAL SKILLS-----------
\section{Технические навыки}
\begin{itemize}[leftmargin=0pt, label={}]
  \small{\item{

          \textbf{Языки и инструменты:}
          Python (средний) --- ООП, ETL, API, обработка данных;
          SQL (средний) --- JOIN, GROUP BY, CTE, оконные функции, query plan
          \\

          \textbf{ETL/ELT:}
          Airflow, dbt
          \\

          \textbf{Хранилища и СУБД:}
          PostgreSQL, ClickHouse
          \\

          \textbf{Big Data:}
          Spark (PySpark), Hadoop
          \\

          \textbf{DevOps и инфраструктура:}
          Git (средний) --- branching, remote repos;
          CI/CD, Docker (продвинутый), Linux/Unix (средний) --- CLI, firewall, SSH, security
          \\

          \textbf{BI:}
          Superset
          \\

        }}
\end{itemize}

%-----------EDUCATION & COURSES-----------
\section{Образование}
\resumeSubHeadingListStart

\resumeSubheading
{Российский государственный аграрный университет}{сен 2022 - июнь 2026}
{Бакалавр, информационные системы и технологии}{}

\resumeSubheading
{Яндекс.Практикум <<Специалист по Data Science>>}{фев 2023 - сен 2023}
{Курс по ML и аналитике данных, \resumelink{https://github.com/gndvrn/yandex-practicum}{репозиторий с проектами}}{}

\resumeSubHeadingListEnd

%-----------ADDITIONAL INFORMATION-----------
\section{Доп. информация}
\resumeItemListStart
\resumeItem{Английский --- upper-intermediate (свободное чтение тех. документации/книг)}
\resumeItem{Понимаю принципы проектирования DWH: Inmon/Kimball, data lake/lakehouse/vault, dimensional modelling, SCD, Medallion, batch/streaming}
\resumeItemListEnd

%-------------------------------------------
\end{document}
